
\section{Appendix to Section \ref{sec:tapediagrams}}\label{app:sec:tapediagrams}


\begin{proof}[Proof of Proposition~\ref{prop:quotienttheory}]
Let $\CatT{\Diag{\Sigma}}_{\tapeeqT}$ be the quotient of $\CatT{\Diag{\Sigma}}$ by $\tapeeqT$: objects are those of  $\CatT{\Diag{\Sigma}}$, arrows are $\tapeeqT$-equivalence classes of arrows of $\CatT{\Diag{\Sigma}}$, hereafter denoted by $[\t]_{\tapeeqT}$. We prove that $\CatT{\Diag{\Sigma}}_{\tapeeqT}$ is isomorphic to $\CatT{\Diag{\mathbb{T}}}$.

We first define the functor $F \colon \CatT{\Diag{\Sigma}}_{\tapeeqT} \to \CatT{\Diag{\mathbb{T}}}$ inductively as
\[\begin{array}{ccccc}
F(\diagp{U})\defeq \; \diagp{U} &  F(\,\bangp{U})\defeq \bangp{U} & F(\tapeFunct{c})\defeq \tapeFunct{[c]_{\mathbb{T}}} &  F(\cobang{U})\defeq \cobang{U} &
F(\codiag{U})\defeq\codiag{U}  \\
 F(\id{U})\defeq \id{U} & F(\id{\zero})\defeq\id{\zero} & F(\sigma_{U,V}^{\piu})\defeq \sigma_{U,V}^{\piu}& F(\t_1 ; \t_2)\defeq F(\t_1);F(\t_2) & F(\t_1 \piu \t_2) \defeq F(\t_1)\piu F(\t_2)
\end{array}\]
Since arrows in $\CatT{\Diag{\Sigma}}_{\tapeeqT}$ are $\tapeeqT$-equivalence classes, we need to prove that the functor is well-defined, namely that if $\s \tapeeqT  \t$ then $F(\s)=F(\t)$. This is done by induction on the rules in \eqref{eq:congr2}. For the case of the rule ($\mathbb{T}$), observe that $\s=\tape{c}$, $\t=\tape{d}$ and $c=_{\mathbb{T}}d$. By the definition of $F$, $F(\s)=F(\t)$. All the other cases are trivial, with the only exception of the rule $(\otimes)$ where one observes that, since  $=_{\mathbb{T}}$ is a congruence w.r.t. $\otimes$, then $F$ is a morphism of rig categories.

The functor  $G \colon  \CatT{\Diag{\mathbb{T}}} \to \CatT{\Diag{\Sigma}}_{\tapeeqT} $ is defined on $\tapeFunct{[c]_{\mathbb{T}}}$ as $G(\tapeFunct{[c]_{\mathbb{T}}})\defeq [\tapeFunct{c}]_{\tapeeqT}$ and then inductively in the same style of the functor $F$ above. Observe that $F(G(\tapeFunct{[c]_{\mathbb{T}}}))=\tapeFunct{[c]_{\mathbb{T}}}$ and $G(F([\tapeFunct{c}]_{\tapeeqT})) = [\tapeFunct{c}]_{\tapeeqT}$. A simple inductive argument confirms that $F$ and $G$ are inverses of each other.
\end{proof}